% Chapter 10 -- The Range Commands (AN-08, AN-09).
% Sources (Law 1): os/i123/app/face_a.c menu_select / range_accept (every
% prompt text quoted below; the format codes of /Range Format and their
% decimal-place prompt 0..15; Date 1-3 and Reset; the label prefixes of
% /Range Label; Name Create / Delete / Labels / Reset / Table; Prot and
% Unprot setting the cell protection bit; Value and Trans through
% i123_transfer; the notices "Illegal number", "Invalid range",
% "Invalid range name", "Command refused", "Not enough worksheet memory"),
% os/i123/core/fmt.c (how each format displays a number: i123_format_number,
% asterisks when the number does not fit, the +/- bar), os/i123/core/
% range.c (names: i123_name_set / delete / labels / table; the name
% retargeting rule), spec/i123/era_shortfalls.md I123-P3-04 (Justify,
% Input, Search refuse), I123-P3-05 (international Date 4/5 and Time
% refuse), I123-P3-06 (range names).  Menu tree generated from the shipped
% resource.
% GATES: test-i123-range-menu test-i123-wk1-roundtrip test-i123-copy-move-model test-i123-menu-tree test-pubs-i123 test-pubs-i123-examples test-pubs-screens
\Gates{test-i123-range-menu, test-i123-wk1-roundtrip, test-i123-copy-move-model,
  test-i123-menu-tree, test-pubs-i123, test-pubs-i123-examples, test-pubs-screens}
\chapter[The Range Commands]{\textit{The Range Commands}}
\input{i123-tree-Range}

The Range commands work on a cell or a rectangle of cells: they set how
numbers and labels are displayed, erase entries, name ranges, protect cells
from change, and copy the values of formulas.

\begin{LotTable}{p{30mm}p{98mm}}
\LotHead{\textbf{Command} & \textbf{Task}}
/Range Format & Sets how the numbers in a range are displayed\\
/Range Label & Sets the alignment of the labels in a range\\
/Range Erase & Erases the entries in a range\\
/Range Name & Names ranges, deletes names and lists them\\
/Range Justify & Rearranges a column of labels \draft{I123-P3-04: refuses}\\
/Range Prot & Marks a range as protected\\
/Range Unprot & Marks a range as unprotected\\
/Range Input & Restricts entry to unprotected cells \draft{I123-P3-04: refuses}\\
/Range Value & Copies a range, replacing formulas by their values\\
/Range Trans & Copies a range, exchanging its rows and columns\\
/Range Search & Opens the menu for finding text in a range \draft{I123-P3-04: its commands refuse}\LotEnd
\end{LotTable}

\section{Common Uses of the Range Commands}
\begin{itemize}
\item Displaying amounts of money with dollar signs and two decimal places
  (\slashcmd{/Range Format Currency})
\item Showing dates as dates rather than serial numbers
  (\slashcmd{/Range Format Date})
\item Centering column headings over the numbers (\slashcmd{/Range Label
  Center})
\item Clearing a block of entries (\slashcmd{/Range Erase})
\item Naming a range so that formulas are easier to read
  (\slashcmd{/Range Name Create})
\item Keeping formulas from being typed over (\slashcmd{/Range Unprot}
  with \slashcmd{/Worksheet Global Protection Enable})
\item Freezing the results of formulas as numbers (\slashcmd{/Range Value})
\end{itemize}

\section{Related Reading}
\begin{itemize}
\item ``Ranges'' and ``Formulas'' in Chapter~1
\item \slashcmd{/Worksheet Global Format} and \slashcmd{/Worksheet Global
  Protection} in Chapter~12
\item Lessons~9, 13, 16, 18 to 21 and 27 in the \doctitle{Tutorial}
\end{itemize}

\section{Terms You Need to Know}
\begin{itemize}
\item A \term{format} is the way a cell displays a number: with a fixed
  number of decimal places, with a dollar sign, as a percentage, as a date.
  A format changes only the display; the number stored in the cell is not
  rounded.
\item The \term{global format} is the format of every cell that has not been
  given one with /Range Format.  It is set with
  \slashcmd{/Worksheet Global Format}.
\item A \term{range name} is a name of up to 15 characters that stands for
  a range in prompts and formulas.
\item A \term{protected cell} cannot be changed while worksheet protection is
  enabled.  Every cell is protected until you unprotect it; protection has
  no effect until you enable it with
  \slashcmd{/Worksheet Global Protection Enable}.
\end{itemize}

\section{Before You Use the Range Commands}
Each Range command asks for the range it is to work on
(Figure~\ref{fig:rangeprompt}).  The control panel shows the address of the
cell pointer: press \kname{ENTER} to work on that one cell, type a range,
such as \type{B1..B3}, or a range name, and press \kname{ENTER}, or move the
cell pointer to another cell with the arrow keys and press \kname{ENTER}.
\FlairShot[fig:rangeprompt]{rangeprompt}{The range prompt of /Range Erase}

\MenuEntry{/Range}
/Range displays the Range menu.

\MenuEntry{/Range Format}
/Range Format sets the format of the cells in a range.
\FlairShot[fig:currency]{currency}{B1..B3 in Currency format with two
decimal places}

\BHead{Terms You Need to Know}
\begin{itemize}
\item The number of \term{decimal places} is the number of digits shown
  after the decimal point, from 0 to 15.
\end{itemize}
A format applies to the cell, not to its present entry: a number you type
into the cell later is displayed in the same format.  A cell keeps its
protection setting when you format it.  A number too wide for its column in
its format is displayed as asterisks; widen the column with
\slashcmd{/Worksheet Column Set-Width} to see it.

\BHead{Procedure}
\begin{Steps123}
\item Select /Range Format and then the format.
\item For Fixed, Sci, Currency, the comma format and Percent, Initech 123
  asks \texttt{Enter number of decimal places (0..15):}.  Type a number from
  0 to 15 and press \kname{ENTER}.
\item At \texttt{Enter range to format:} specify the range and press
  \kname{ENTER}.
\end{Steps123}

\begin{Options}
\Option{Fixed} A fixed number of decimal places: \texttt{-1234.57}.
\Option{Sci} Scientific notation, a number between 1 and 10 times a power
  of ten: \texttt{-1.23E+03}.
\Option{Currency} A dollar sign, commas between thousands and a fixed
  number of decimal places; negative numbers in parentheses:
  \texttt{(\$1,234.57)}.
\Option{,} Commas between thousands and a fixed number of decimal places;
  negative numbers in parentheses: \texttt{(1,234.57)}.
\Option{General} As many decimal places as the number has and the column
  allows; scientific notation for numbers too large or too small:
  \texttt{-1234.567}.
\Option{+/-} A bar of plus signs, one for each whole unit of a positive
  number, or of minus signs for a negative number, rounded to the nearest
  whole unit; a period when the number rounds to zero.
\Option{Percent} The number times 100, with a percent sign and a fixed
  number of decimal places: \texttt{45.7\%}.
\OptionMenu{Date} Displays the Date menu (below).
\Option{Text} A formula instead of its value; a number in General format.
\Option{Hidden} Nothing: the cell appears empty, though its entry is still
  there and is shown on the control panel.
\Option{Reset} Returns the cells to the global format.
\end{Options}

The same number in each format, in a column 12 characters wide (the cell
contents, then what the cell shows):
\LotExample{examples/formats}

\SubEntry{/Range Format Date}
Initech 123 stores a date as a \term{date serial number}, the number of days
since 31 December 1899: 1 January 1900 is 1.  A date format displays a
serial number as a date.  Use \fn{@DATE} to enter a date.
\begin{Options}
\Option{1 (DD-MMM-YY)} Day, month and year: \texttt{02-Mar-92}.  The column
  must be at least 10 characters wide.
\Option{2 (DD-MMM)} Day and month: \texttt{02-Mar}.
\Option{3 (MMM-YY)} Month and year: \texttt{Mar-92}.
\Option[I123-P3-05: refuses]{4 (Long Intn'l)} The international long date
  format.
\Option[I123-P3-05: refuses]{5 (Short Intn'l)} The international short
  date format.
\OptionMenu{Time} Displays the Time menu (below).
\end{Options}
\LotExample{examples/datefmt}

\SubEntry{/Range Format Date Time}
\begin{Options}
\Option[I123-P3-05: refuses]{1 (HH:MM:SS AM/PM)} Hours, minutes and seconds
  in 12-hour form.
\Option[I123-P3-05: refuses]{2 (HH:MM AM/PM)} Hours and minutes in 12-hour
  form.
\Option[I123-P3-05: refuses]{3 (Long Intn'l)} The international long time
  format.
\Option[I123-P3-05: refuses]{4 (Short Intn'l)} The international short time
  format.
\end{Options}

\LabelHead{Tip}
\begin{itemize}
\item The +/- format makes a rough bar chart of small whole numbers:
\end{itemize}
\LotExample{examples/bargraph}

\MenuEntry{/Range Label}
/Range Label changes the label prefix, and so the alignment, of every label
in a range.  Numbers and formulas in the range are not affected, and labels
you type in the range later get the default prefix.
\FlairShot[fig:labelmenu]{labelmenu}{The /Range Label menu}
\BHead{Procedure}
\begin{Steps123}
\item Select /Range Label and then Left, Right or Center.
\item At \texttt{Enter range of labels:} specify the range and press
  \kname{ENTER}.
\end{Steps123}
\begin{Options}
\Option{Left} Gives the labels the prefix \texttt{'}: aligned at the left.
\Option{Right} Gives the labels the prefix \texttt{"}: aligned at the right.
\Option{Center} Gives the labels the prefix \texttt{\textasciicircum}:
  centered.
\end{Options}

\section{Aligning Headings: An Example}
Three headings typed without a prefix are left-aligned.  /Range Label
Center gives all three the caret prefix; /Range Label Right then gives the
two over the numbers the quotation-mark prefix:

\LotExample{examples/labelalign}

\MenuEntry{/Range Erase}
/Range Erase erases the entries in a range.  The cells keep their formats
and protection settings.
\BHead{Procedure}
\begin{Steps123}
\item Select /Range Erase.
\item At \texttt{Enter range to erase:} specify the range and press
  \kname{ENTER}.
\end{Steps123}
\begin{Caution}
Formulas that refer to an erased cell treat it as empty.  Save the
worksheet before you erase entries you may want back.
\end{Caution}

\MenuEntry{/Range Name}
/Range Name gives names to ranges, deletes them and lists them
(Figure~\ref{fig:namemenu}).
\FlairShot[fig:namemenu]{namemenu}{The /Range Name menu}
\FlairShot[fig:nametable]{nametable}{The range name COSTS and its table}

\BHead{Terms You Need to Know}
\begin{itemize}
\item A \term{range name} has up to 15 characters.  Initech 123 stores it
  in capital letters; you can type it in either.
\end{itemize}
A formula can use a range name wherever it can use a range or a cell
address: \type{@SUM(COSTS)}.  The control panel shows such a formula with
the range the name stands for.  A name typed with a dollar sign before it,
\type{\$COSTS}, is an absolute reference to the range.

\begin{Options}
\Option{Create} Names a range, or changes the range of an existing name.
  At \texttt{Enter name:} type the name and press \kname{ENTER}; at
  \texttt{Enter range:} specify the range and press \kname{ENTER}.
\Option{Delete} Deletes one name.  At \texttt{Enter name to delete:} type
  the name and press \kname{ENTER}.  Formulas that used the name keep the
  range it stood for.
\OptionMenu{Labels} Displays the Labels menu (below).
\Option{Reset} Deletes every range name at once.  Formulas keep the ranges
  the names stood for.
\Option{Table} Lists the range names in the worksheet.  At
  \texttt{Enter range for table:} specify the upper left cell; Initech 123
  enters each name in that column and its range in the column to the
  right.
\end{Options}

\SubEntry{/Range Name Labels}
/Range Name Labels names single cells after labels next to them: each
label in the range you specify becomes the name of the cell beside it.  At
\texttt{Enter label range:} specify the range of labels.
\begin{Options}
\Option{Right} Each label names the cell to its right.
\Option{Down} Each label names the cell below it.
\Option{Left} Each label names the cell to its left.
\Option{Up} Each label names the cell above it.
\end{Options}

\begin{Note}
A name that is longer than 15 characters displays \texttt{Invalid range
name}.  While the worksheet contains formulas, you cannot move an existing
name to a different range with Create: Initech 123 displays
\Msg{Command refused}.  Delete the name and create it again; formulas
that used it keep its old range.
\end{Note}

\section{Naming a Range: An Example}
B1 holds 120 and B2 holds 950.  Name B1..B2 \texttt{COSTS} with
\slashcmd{/Range Name Create}, enter \type{@SUM(COSTS)} in B4, then move the
cell pointer to D1 and select \slashcmd{/Range Name Table} and press
\kname{ENTER}:
\LotExample{examples/names}

\MenuEntry[I123-P3-04: refuses]{/Range Justify}
Rearranges a column of labels so that the text fills lines of a given
width.

\MenuEntry{/Range Prot}
/Range Prot marks the cells of a range as protected again after
\slashcmd{/Range Unprot}.  At \texttt{Enter range to protect:} specify the
range and press \kname{ENTER}.

\MenuEntry{/Range Unprot}
/Range Unprot marks the cells of a range as unprotected, so that you can
change them while worksheet protection is enabled.  At
\texttt{Enter range to unprotect:} specify the range and press
\kname{ENTER}.

\begin{Note}
While protection is enabled, typing an entry into a protected cell displays
\Msg{Cell is protected}, and a command that would change a protected cell
displays \Msg{Command refused}; nothing is changed.
\end{Note}

\MenuEntry[I123-P3-04: refuses]{/Range Input}
Restricts the cell pointer to the unprotected cells of a range, for filling
in a form.

\MenuEntry{/Range Value}
/Range Value copies a range, as \slashcmd{/Copy} does, except that each
formula is copied as its current value: a number, or a label if the formula
produces a string.
\BHead{Procedure}
\begin{Steps123}
\item Select /Range Value.
\item At \texttt{Enter range to copy FROM:} specify the range and press
  \kname{ENTER}.
\item At \texttt{Enter range to copy TO:} specify the upper left cell of the
  destination and press \kname{ENTER}.  The prompt offers the cell the
  cell pointer is in (Figure~\ref{fig:valueto}).
\end{Steps123}
\FlairShot[fig:valueto]{valueto}{The TO prompt of /Range Value}
\LabelHead{Tip}
\begin{itemize}
\item To freeze a block of results, give the same cell for FROM and TO: the
  formulas are replaced by their values.
\end{itemize}

\section{Freezing Results: An Example}
C1 and C2 compute a cost from the quantities in column B.  /Range Value
with C1..C2 as both the FROM and the TO range replaces each formula by its
value: the cells show the same numbers, but the control panel now shows
numbers, not formulas, and a change in column B no longer affects them.

\LotExample{examples/freeze}

\MenuEntry{/Range Trans}
/Range Trans copies a range with its rows and columns exchanged: the first
row of the FROM range becomes the first column of the copy.  Formulas are
copied as their current values.  The prompts are those of /Range Value.

\section{Transposing a Table: An Example}
A1..C2 holds three month names over three values, the third a formula.
/Range Trans with A1..C2 as the FROM range and A4 as the TO range writes
the months down column A from A4 and the values beside them; the formula
in C2 arrives as its value, 30:

\LotExample{examples/transpose}

\MenuEntry{/Range Search}
/Range Search displays the menu for finding or replacing text in a range.
\begin{Options}
\OptionMenu{Formulas} Displays the menu for searching formulas (below).
\Option[I123-P3-04: refuses]{Label} Searches the labels only.
\Option[I123-P3-04: refuses]{Both} Searches formulas and labels.
\end{Options}
\SubEntry{/Range Search Formulas}
\begin{Options}
\Option[I123-P3-04: refuses]{Find} Highlights each occurrence of the search
  text.
\Option[I123-P3-04: refuses]{Replace} Replaces each occurrence of the search
  text.
\end{Options}
